\chapter{Минимизация логических функций}
\label{ch:minimize}

В главе~\ref{ch:comb} мы строили схемы по таблице истинности и упрощали их по
карте Карно для трёх переменных. На практике функций от четырёх и более
переменных гораздо больше, а у некоторых входных комбинаций вообще нет
определённого значения. Здесь мы разберём полный набор приёмов: четыре
переменные, безразличные состояния, формальный алгоритм Квайна---Мак-Класки и
то, как всё это связано с тем, что делает синтезатор в ПЛИС.

\section{Канонические формы}

Любую функцию можно записать \emph{в лоб}, прямо по таблице истинности.

\begin{itemize}
  \item \term{Совершенная дизъюнктивная нормальная форма} (СДНФ): для каждой
        строки, где $F=1$, пишем произведение всех переменных (с инверсиями у
        тех, что равны нулю), и складываем эти произведения.
        Каждое такое произведение --- \term{минтерм} $m_i$.
  \item \term{Совершенная конъюнктивная нормальная форма} (СКНФ): для каждой
        строки, где $F=0$, пишем сумму всех переменных (с инверсиями у тех,
        что равны \emph{единице}), и перемножаем эти суммы. Каждая такая сумма
        --- \term{макстерм} $M_i$.
\end{itemize}

Эти формы всегда правильны, но избыточны: у функции от четырёх переменных
СДНФ может состоять из восьми слагаемых по четыре буквы в каждом. Задача
минимизации --- найти запись с наименьшим числом слагаемых и букв.
Функцию часто задают списком номеров строк: запись $F=\Sigma m(0,1,2,5)$
означает «функция равна единице на строках 0, 1, 2 и 5, на остальных нулю».

\section{Карта Карно на четыре переменные}

Для четырёх переменных $A,B,C,D$ карта состоит из 16 клеток в виде квадрата
$4\times 4$. Строки подписаны значениями $AB$, столбцы --- $CD$, и те и
другие в порядке кода Грея: \code{00, 01, 11, 10}. Благодаря этому клетки,
соседние по горизонтали и вертикали, отличаются одной переменной. Важное
следствие: \emph{карта замкнута}. Левый столбец соседствует с правым, верхняя
строка --- с нижней, а четыре угла вместе образуют одну группу.

Минимизируем функцию $F = \Sigma m(0,1,2,5,8,9,10)$.

\begin{figure}[H]
\centering
\begin{tikzpicture}[cell/.style={draw, minimum size=10mm, font=\sffamily\large}]
  % подписи
  \foreach \c/\lab in {0/00, 1/01, 2/11, 3/10} \node[font=\sffamily\small] at (\c,0.85) {\lab};
  \foreach \r/\lab in {0/00, 1/01, 2/11, 3/10} \node[font=\sffamily\small] at (-0.95,-\r) {\lab};
  \node[font=\sffamily\small] at (-1.0,0.85) {$AB\backslash CD$};
  % значения F
  \foreach \r/\c/\v in {0/0/1, 0/1/1, 0/2/0, 0/3/1,
                        1/0/0, 1/1/1, 1/2/0, 1/3/0,
                        2/0/0, 2/1/0, 2/2/0, 2/3/0,
                        3/0/1, 3/1/1, 3/2/0, 3/3/1}
    {\ifnum\v=1 \def\cc{FpgaLight}\else\def\cc{white}\fi
     \node[cell, fill=\cc] at (\c,-\r) {\v};}
  % группа B'D': четыре угла
  \foreach \x/\y in {0/0, 3/0, 0/-3, 3/-3}
    \draw[very thick, FpgaRed, rounded corners=4pt] (\x-0.44,\y+0.44) rectangle (\x+0.44,\y-0.44);
  % группа B'C': 2 клетки сверху и 2 снизу (замыкание по вертикали)
  \draw[very thick, FpgaGreen, rounded corners=4pt] (-0.40,0.40) rectangle (1.40,-0.40);
  \draw[very thick, FpgaGreen, rounded corners=4pt] (-0.40,-2.60) rectangle (1.40,-3.40);
  % группа A'C'D: два соседа по вертикали
  \draw[very thick, FpgaOrange, rounded corners=4pt] (0.68,0.30) rectangle (1.32,-1.30);
  % пояснения
  \node[font=\sffamily\small, text=FpgaRed, anchor=west] at (4.4,0.5) {$\overline{B}\,\overline{D}$ --- четыре угла};
  \node[font=\sffamily\small, text=FpgaGreen, anchor=west] at (4.4,-0.3) {$\overline{B}\,\overline{C}$ --- верх и низ, два столбца};
  \node[font=\sffamily\small, text=FpgaOrange, anchor=west] at (4.4,-1.1) {$\overline{A}\,\overline{C}\,D$ --- клетки 1 и 5};
  \node[font=\sffamily\small\bfseries, anchor=west] at (4.4,-2.1) {$F=\overline{B}\,\overline{D}+\overline{B}\,\overline{C}+\overline{A}\,\overline{C}D$};
\end{tikzpicture}
\caption{Карта Карно четырёх переменных. Группы красного цвета (четыре угла) и зелёного (две клетки сверху и две снизу) получаются замыканием краёв карты}
\label{fig:karnaugh4}
\end{figure}

Порядок действий такой же, как и для трёх переменных:

\begin{enumerate}
  \item Отметить единицы.
  \item Объединить их в прямоугольные группы по $1, 2, 4, 8$ или 16 клеток;
        группы могут пересекаться и «переползать» через край.
  \item Взять как можно более крупные группы и покрыть все единицы наименьшим
        числом групп.
  \item Каждой группе соответствует произведение только тех переменных,
        которые в группе не меняются. Группа из $2^k$ клеток исключает $k$
        переменных.
\end{enumerate}

Исходная СДНФ состояла из семи слагаемых по четыре буквы (28 букв), а
результат --- из трёх слагаемых по 2--3 буквы (7 букв).

\section{Безразличные состояния}

Бывает, что некоторые комбинации входов вообще не встречаются. Например,
цифра в двоично-десятичном коде не может быть больше девяти: комбинации
\code{1010}--\code{1111} в BCD запрещены. Что на них должен выдавать
дешифратор или компаратор --- всё равно. Такие комбинации называют
\term{безразличными состояниями} (don't care, обозначение $d$ или $\times$).
Их можно считать как нулями, так и единицами --- в зависимости от того, что
позволяет получить более крупные группы.

\begin{figure}[H]
\centering
\begin{tikzpicture}[cell/.style={draw, minimum size=10mm, font=\sffamily\large}]
  \foreach \c/\lab in {0/00, 1/01, 2/11, 3/10} \node[font=\sffamily\small] at (\c,0.85) {\lab};
  \foreach \r/\lab in {0/00, 1/01, 2/11, 3/10} \node[font=\sffamily\small] at (-0.95,-\r) {\lab};
  \node[font=\sffamily\small] at (-1.0,0.85) {$AB\backslash CD$};
  % «цифра больше четырёх»: единицы 5..9, безразличные 10..15
  \foreach \r/\c/\v in {0/0/0, 0/1/0, 0/2/0, 0/3/0,
                        1/0/0, 1/1/1, 1/2/1, 1/3/1,
                        2/0/x, 2/1/x, 2/2/x, 2/3/x,
                        3/0/1, 3/1/1, 3/2/x, 3/3/x}
    {\def\cc{white}
     \if 1\v \def\cc{FpgaLight}\fi
     \if x\v \def\cc{gray!25}\fi
     \node[cell, fill=\cc] at (\c,-\r) {\if x\v$\times$\else\v\fi};}
  \draw[very thick, FpgaRed, rounded corners=4pt] (-0.40,-1.60) rectangle (3.40,-3.40);
  \draw[very thick, FpgaGreen, rounded corners=4pt] (1.60,-0.60) rectangle (3.40,-2.40);
  \draw[very thick, FpgaOrange, rounded corners=4pt] (0.60,-0.62) rectangle (2.40,-2.38);
  \node[font=\sffamily\small, text=FpgaRed, anchor=west] at (4.4,-0.3) {$A$ --- нижние строки (8--15)};
  \node[font=\sffamily\small, text=FpgaGreen, anchor=west] at (4.4,-1.0) {$BC$ --- столбцы $CD=11,10$};
  \node[font=\sffamily\small, text=FpgaOrange, anchor=west] at (4.4,-1.7) {$BD$ --- столбцы $CD=01,11$};
  \node[font=\sffamily\small\bfseries, anchor=west] at (4.4,-2.5) {$F = A + BC + BD$};
\end{tikzpicture}
\caption{Функция «цифра BCD больше четырёх». Серые клетки --- безразличные состояния $10,\ldots,15$. Благодаря им группы получаются крупными, и результат состоит всего из трёх слагаемых}
\label{fig:karnaugh_dc}
\end{figure}

Без безразличных состояний пришлось бы писать пять отдельных
минтермов, а с ними функция превратилась в $F = A + BC + BD$. Поэтому при
описании схем разработчик всегда должен спрашивать себя: «все ли входные
комбинации возможны?». В Verilog безразличное состояние записывают как
\code{x}: например, \code{default: y = 4'bx;} разрешает синтезатору выбрать
любое значение, и схема получается меньше.

\section{Метод Квайна---Мак-Класки}

Карта Карно удобна для четырёх, ну пяти переменных. Для десяти она
превращается в мучение, а для двадцати бесполезна. Поэтому синтезаторы
используют формальный метод, который можно выполнить на компьютере:
\term{метод Квайна---Мак-Класки}. Он состоит из двух частей.

\textbf{Часть 1. Поиск простых импликант.} Запишем минтермы двоичными
числами и сгруппируем по числу единиц. Склеим любые два числа, которые
отличаются ровно в одном разряде: этот разряд заменяется чертой \code{-}
(«переменная не важна»). Повторяем склейку, пока она возможна. Всё, что не
склеилось, --- \term{простая импликанта}.

\begin{table}[H]
\centering
\caption{Метод Квайна---Мак-Класки для $F=\Sigma m(0,1,2,5,8,9,10)$}
\small
\begin{tabularx}{\textwidth}{l l l l}
\toprule
\textbf{Единиц} & \textbf{Минтермы} & \textbf{Склейка 1} & \textbf{Склейка 2} \\
\midrule
0 & 0: \code{0000} & 0,1: \code{000-} & 0,1,8,9: \code{-00-} \\
  &               & 0,2: \code{00-0} & 0,2,8,10: \code{-0-0} \\
  &               & 0,8: \code{-000} & \\
\midrule
1 & 1: \code{0001} & 1,5: \code{0-01} (простая) & \\
  & 2: \code{0010} & 1,9: \code{-001} & \\
  & 8: \code{1000} & 2,10: \code{-010} & \\
  &               & 8,9: \code{100-} & \\
  &               & 8,10: \code{10-0} & \\
\midrule
2 & 5: \code{0101} & & \\
  & 9: \code{1001} & & \\
  & 10: \code{1010} & & \\
\bottomrule
\end{tabularx}
\end{table}

Простые импликанты: \code{-00-} ($\overline{B}\,\overline{C}$), \code{-0-0}
($\overline{B}\,\overline{D}$) и \code{0-01} ($\overline{A}\,\overline{C}D$)
--- те же три группы, которые мы нашли на карте.

\textbf{Часть 2. Выбор покрытия.} Строим таблицу: строки --- простые
импликанты, столбцы --- минтермы. Если какой-то минтерм покрывается только
одной импликантой, она \term{существенная} и обязательно входит в ответ.
Здесь минтерм 5 покрывает только \code{0-01}, минтермы 2 и 10 --- только
\code{-0-0}, а 1, 9 --- только \code{-00-}. Значит, все три импликанты
существенны, и минимальная форма совпала с результатом по карте.

Метод находит оптимальный ответ, но при росте числа переменных число
импликант растёт лавинообразно. Поэтому в реальных программах используют
эвристики, которые не гарантируют абсолютный минимум, но работают очень быстро
на функциях с сотнями входов (Espresso, ABC).

\section{Многоуровневые схемы}

Минимальная форма в виде суммы произведений даёт схему в \emph{два уровня}:
сначала элементы И, потом ИЛИ. Она быстрая, но может требовать элементов с
очень большим числом входов. Если вынести общий множитель за скобки, схема
становится многоуровневой, зато экономнее.

\begin{figure}[H]
\centering
\begin{tikzpicture}[font=\sffamily\small, x=1cm, y=1cm]
  \node at (0,2.2) {\textbf{Два уровня}};
  \node at (0,1.6) {$F = AB + AC + AD$};
  \node[draw, rounded corners=2pt, minimum width=14mm] at (-1.5,0.5) {И};
  \node[draw, rounded corners=2pt, minimum width=14mm] at (0,0.5) {И};
  \node[draw, rounded corners=2pt, minimum width=14mm] at (1.5,0.5) {И};
  \node[draw, rounded corners=2pt, minimum width=22mm] at (0,-0.8) {ИЛИ (3 входа)};
  \draw (-1.5,0.2) -- (-1.5,-0.5); \draw (0,0.2) -- (0,-0.5); \draw (1.5,0.2) -- (1.5,-0.5);
  \node at (0,-1.5) {4 элемента, 9 входов};
  \begin{scope}[xshift=7cm]
  \node at (0,2.2) {\textbf{Три уровня}};
  \node at (0,1.6) {$F = A\,(B + C + D)$};
  \node[draw, rounded corners=2pt, minimum width=22mm] at (0,0.5) {ИЛИ (3 входа)};
  \node[draw, rounded corners=2pt, minimum width=14mm] at (0,-0.8) {И};
  \draw (0,0.2) -- (0,-0.5);
  \node at (0,-1.5) {2 элемента, 5 входов};
  \end{scope}
\end{tikzpicture}
\caption{Вынесение общего множителя уменьшает число элементов, но добавляет ещё один этап задержки}
\end{figure}

Здесь стоит помнить компромисс: чем глубже схема, тем она меньше по площади,
но медленнее по задержке. Поэтому синтезатор ищет золотую середину в зависимости
от требуемой тактовой частоты.

\begin{infobox}[Как это выглядит в ПЛИС]
В ПЛИС «элемент» --- это LUT с четырьмя входами, которая умеет реализовать
\emph{любую} функцию четырёх переменных за одну ячейку. Поэтому важно не число
слагаемых в минимальной форме, а число переменных, от которых зависит
функция: нужно уложить логику в минимальное число LUT4. Если функция зависит
от шести переменных, её придётся разложить на две или три LUT. Эту работу
синтезатор Yosys выполняет программой ABC.
\end{infobox}

\section{Вопросы и задачи}
\label{sec:tasks-minimize}

Ответы собраны в приложении~\ref{app:answers}.

\begin{taskbox}[Задачи к главе]
\begin{enumerate}
  \item Минимизируйте $F_1=\Sigma m(0,2,8,10)$, $F_2=\Sigma m(1,3,5,7,9,11,13,15)$ и $F_3=\Sigma m(0,1,2,3)$ от четырёх переменных $A,B,C,D$.
  \item Функция «цифра BCD чётна» задана минтермами $0,2,4,6,8$, а комбинации $10$--$15$ безразличны. Найдите минимальную форму.
  \item Найдите простые импликанты функции $F=\Sigma m(0,1,2,5,6,7)$ от трёх переменных методом Квайна---Мак-Класки и подберите минимальное покрытие.
  \item Приведите пример функции, для которой по карте Карно образуется группа, замыкающаяся через край карты.
  \item Почему для ПЛИС важнее число переменных функции, чем число слагаемых в минимальной форме?
\end{enumerate}
\end{taskbox}
